% Generated draft: requires mathematical and source review.
\documentclass[11pt]{article}
\usepackage[margin=25mm]{geometry}
\usepackage{fontspec}
\setmainfont{Latin Modern Roman}
\usepackage{amsmath,amssymb,mathtools,unicode-math}
\setmathfont{Latin Modern Math}
\usepackage{enumitem,needspace,xurl,hyperref}
\setlength{\parindent}{0pt}
\setlength{\parskip}{5pt}
\setlength{\emergencystretch}{3em}
\newcommand{\N}{\mathbb{N}}
\newcommand{\Z}{\mathbb{Z}}
\newcommand{\Q}{\mathbb{Q}}
\newcommand{\R}{\mathbb{R}}
\newcommand{\C}{\mathbb{C}}
\newcommand{\F}{\mathbb{F}}
\newcommand{\A}{\mathbb{A}}
\newcommand{\PP}{\mathbb{P}}
\newcommand{\E}{\mathbb{E}}
\newcommand{\eps}{\varepsilon}
\newcommand{\dd}{\,\mathrm{d}}
\begin{document}
\section*{UnsolvedMath ID: 800986}
\textbf{Source problem code:} catalogue entry
\subsection{Definitions and mathematical statement}
The supplied source asks: \emph{In dimension \(2\), is each irreversible path structure metrizable?} The meanings of ``path structure'' and ``metrizable'' are not defined in the supplied excerpt, and the question itself does not explicitly repeat the locality qualification in the preceding sentence.

A precise conditional formulation under the smooth, local Finsler interpretation is as follows. Let \(M\) be a smooth surface. An oriented path structure \(\mathcal P\) is a smoothly varying family of oriented, unparameterized immersed curves such that each point and each positive tangent ray determine a unique local path germ. Here a positive tangent ray is an equivalence class of nonzero tangent vectors under multiplication by positive real numbers. The structure is reversible if reversing the orientation of every path preserves the family of path germs, and irreversible otherwise.

Under this interpretation, local Finsler metrizability means that, near every point \(p\in M\), there is a neighborhood \(U\) and a function \(F:TU\to[0,\infty)\), continuous everywhere and smooth away from the zero section, satisfying
\[
F(x,0)=0,\qquad F(x,v)>0\quad(v\ne0),\qquad F(x,\lambda v)=\lambda F(x,v)\quad(\lambda>0),
\]
with positive definite fundamental tensor
\[
g_{ij}(x,v)=\frac12\frac{\partial^2(F^2)}{\partial v^i\partial v^j}(x,v)\qquad(v\ne0),
\]
such that the oriented, unparameterized geodesic germs of \(F\) are exactly the path germs of \(\mathcal P\) in \(U\). No condition \(F(x,v)=F(x,-v)\) is imposed.

The corresponding conditional question is: for every smooth surface \(M\), every irreversible oriented path structure \(\mathcal P\) on \(M\), and every point \(p\in M\), does such a neighborhood \(U\) and such an \(F\) exist?

The supplied excerpt does not establish that this is its intended definition of metrizability, its intended regularity class, or its intended local quantification. It also specifies no boundary conditions. These qualifications require confirmation from the full source. The excerpt reports that every reversible path structure in dimension \(2\) is locally metrizable, citing [Alv-Ber2010, Theorem B]; this is retained as a supplied background claim.
\subsection{Short English statement}
The question asks whether every irreversible path structure on a surface is metrizable. Under the explicitly conditional local Finsler interpretation, this asks whether its oriented paths can locally be exactly the geodesics of a possibly nonreversible Finsler metric. The intended metric class and the local versus global scope require source confirmation.
\subsection{Sources}
\begin{itemize}
\item Local UnsolvedMath registry, supplied record with permanent ID \(800986\), source problem code \texttt{catalogue entry}, in the AMR Open Problem Lists collection. The supplied status is \texttt{open}; this label is not an independently verified mathematical assessment.
\item Burns and Matveev, \emph{Open Problems and Questions about Geodesics}, attributed in the supplied record to 2019, Question 67. Supplied source URL: \url{https://arxiv.org/abs/1308.5417}. The supplied provenance records access on 2026-10-16 and extraction by \texttt{arxiv\_latex}; neither access nor source contents have been independently verified here.
\item [Alv-Ber2010, Theorem B], cited by the supplied excerpt for local metrizability of reversible path structures in dimension \(2\). No full bibliographic entry or separate URL for this reference is supplied.
\end{itemize}
\end{document}
